\subsection{Adding bottom components to the cell vectors}\label{sec:bottom-shift-obstruction}
For the two examples with the smallest displayed supported-zero counts, a further family of changes of basis can be excluded even when the new bottom line need not be an ideal.

\begin{proposition}\label{prop:bottom-shift-obstruction}
Start with dimension algebra B or C over $\mathbb Q(q)$. Keep the unit, replace each cell vector by
\[
 \gamma\longmapsto\gamma+r_\gamma\Sigma,
 \qquad r_\gamma\in\mathbb Q(q),
\]
and replace $\Sigma$ by any vector completing these nine vectors to a basis. Suppose the resulting structure coefficients are nonnegative Laurent polynomials and give an effective lift of the fixed rank-three trace table. Among the three bottom coefficients of
\[
 (\gamma_1^1)^2,\qquad
 \gamma_1^1\gamma_{12}^1,\qquad
 \gamma_{12}^1\gamma_1^1
\]
in the new basis, at least one has a supported zero above its trace-support interval and at least one has a supported zero below its trace-support interval. In particular there are at least two bottom supported-zero weights, and this class of changes cannot yield exact bottom supports.
\end{proposition}

\begin{proof}
Evaluate at a sufficiently large or sufficiently small positive real $q$, avoiding the finitely many exceptional points of the rational change of basis. All the proposed new structure constants are nonnegative, so the cone generated by the new basis is pointed and closed under multiplication.

In the old algebra put $e=\Sigma/S$, where $\Sigma^2=S\Sigma$. Then $e$ is a nonzero central idempotent and $S(q)>0$. Products between vectors from different old cells lie in $\mathbb Q(q)e$. Adding multiples of $\Sigma$ to the cell vectors preserves this property. One such cross-cell product has a strictly positive new bottom coefficient, because its prescribed trace coefficient is nonzero. Consequently a nonzero scalar multiple of $e$ belongs to the new positive cone. It follows that either $e$ or $-e$ is positive. The latter alternative is impossible: closure under multiplication would also make $e=(-e)^2$ positive, contradicting pointedness. Thus $e$ is positive, and its new bottom coordinate $t$ is strictly positive.

Let $X,Y$ denote the new cell vectors labelled $\gamma_1^1,\gamma_{12}^1$. There are scalars $u_x,u_y$ with
\[
 eX=u_xe,\qquad eY=u_ye.
\]
Since $e,X,Y$ are positive and $e\ne0$, one has $u_x,u_y\ge0$. The quotient by $\mathbb Q(q)e$ retains the relations
\[
 X^2=cX,\qquad XY=cY,\qquad YX=cX,
 \qquad c=dD>0,
\]
where $d=D(D-3)$ in B and $d=(D-1)(D-2)$ in C. These follow directly from $P_1^2=DP_1$, $P_1P_{12}=DP_{12}$, $P_{12}P_1=DP_1$, and $\lambda_1^+=D$.

Write $P,Q,W$ for the new bottom coefficients of $X^2,XY,YX$. Multiplying the three quotient relations by $e$ determines their missing scalar terms. Because $X$ and $Y$ themselves have zero bottom coordinate in the new basis, this gives
\[
 P=t u_x(u_x-c),\qquad
 Q=t u_y(u_x-c),\qquad
 W=t u_x(u_y-c).
\]
All three are strictly positive: their fixed traces are $D,D^2,5D-4$. Hence $u_x>c$ and $u_y>c$. Moreover
\[
 Q-P=t(u_x-c)(u_y-u_x),\qquad
 W-Q=tc(u_y-u_x).
\]
The two differences therefore have the same sign.

The trace supports of $P,Q,W$ are respectively $[-1,1]$, $[-2,2]$, and $[-1,1]$, intersected with $\mathbb Z$. If none of these coefficients had a supported zero above its trace interval, their maximum exponents would be $1,2,1$. Their positive leading coefficients would imply $Q>P$ and $Q>W$ for all sufficiently large positive $q$, contradicting the same-sign identity. Thus at least one coefficient has a supported zero above its trace interval.

Likewise, if none had a supported zero below its trace interval, their minimum exponents would be $-1,-2,-1$. Letting $q\to0^+$ would again give $Q>P$ and $Q>W$, with the same contradiction. The two resulting weights are distinct because one lies above a positive endpoint and the other below a negative endpoint. This proves the lower bound.
\end{proof}

The proposition permits an arbitrary new bottom vector and does not assume that its span is an ideal. It still concerns a restricted family: mixing cell vectors with one another, adding unit components to them, or changing their individual scalar normalizations may escape this argument.
